\subsection{Análisis de tecnologías del Módulo 2}

En esta sección se presenta un análisis,comparacción y selección
tanto del lenguaje de programación así como del módelo de logica difusa a implementar
,tomando en cuenta los requerimientos funcionales y 
el formato de datos (JSON) previamente seleccionado con el cual va a trabajar el Módulo 1.

\subsubsection{Análisis del módelo de logica difusa}

Con base en la evaluación de la \cref{tab:modelos_logica_difusa}, 
se seleccionó el \textbf{Modelo de Mamdani} para el Módulo 2 por 
ser el único capaz de procesar la alta incertidumbre de la fermentación
del kéfir mediante conjuntos difusos cerrados y ponderaciones por 
rangos (como pH de $4.2$ a $4.6$ y temperatura 
de $18\,^{\circ}\text{C}$ a $30\,^{\circ}\text{C}$), 
mapeando directamente las reglas hacia la superficie de control 
con una elevada fiabilidad. En contraste, se descartó 
el modelo Takagi-Sugeno debido a que no utiliza un conjunto difuso 
consecuente y exige conocer de antemano la superficie de control 
mediante datos históricos preexistentes, los cuales no están 
disponibles en este proceso artesanal. Por su parte, 
el modelo de Tsukamoto requiere estrictamente funciones de 
pertenencia abiertas y monótonas en las salidas, por lo que el
uso de los rangos cerrados propios de este proyecto generaría 
una indefinición matemática al arrojar múltiples valores escalares
para una misma coordenada de corte.

Asimismo, con base en el proyecto presentado en la \cref{sec:In2} 
del estado del arte, se demuestra la elevada fiabilidad del modelo 
de Mamdani en este tipo de bioprocesos, al contar con un antecedente 
exitoso de implementación documentado.

\begin{table}[htbp]
    \centering
    \small
    \caption{Modelos de lógica difusa.}
    \label{tab:modelos_logica_difusa}
    \begin{tabular}{|>{\raggedright\arraybackslash}p{2.6cm}|>{\raggedright\arraybackslash}p{3.8cm}|>{\raggedright\arraybackslash}p{4.2cm}|>{\raggedright\arraybackslash}p{3.8cm}|}
        \hline
        \textbf{Criterio} & \textbf{Modelo Mamdani}\cite{Morcillo} & \textbf{Takagi-Sugeno (TSK)} & \textbf{Modelo Tsukamoto} \\ \hline
        \textbf{Tipo de Consecuente} & Conjunto difuso (etiqueta lingüística). & Función matemática lineal o valor constante. & Conjunto difuso con función monótona. \\ \hline
        \textbf{Geometría de Salida} & Soporta funciones de pertenencia cerradas (rangos). & Sin conjunto difuso consecuente (salida escalar). & Exige funciones abiertas (monótonas continuas). \\ \hline
        \textbf{Requisito de Información} & Basado en reglas expertas e intervalos de operación. & Requiere conocer la superficie de control o datos previos. & Evaluación de un solo valor por nivel de corte $\alpha$. \\ \hline
        \textbf{Desfuzificación} & Integración geométrica (Centroide, Bisectriz). & Promedio ponderado (\textit{Weighted Average}). & Promedio ponderado de puntos de corte. \\ \hline
        \textbf{Limitación Principal} & Mayor costo computacional en la desfuzificación con base a las entradas. & Inviable sin datos históricos del proceso. & Ambigüedad $z_1, z_2$ con rangos cerrados. \\ \hline
        \textbf{Aplicación} & Sistemas de monitoreo, diagnóstico e inferencia cualitativa. & Descartado: Control adaptativo con curva conocida. & Descartado: Sistemas simples con salidas abiertas. \\ \hline
    \end{tabular}
\end{table}


\subsubsection{Análisis del lenguaje de programación}




Con base en la evaluación presentada en la \cref{tab:lenguajesprogramacion}, 
se seleccionó \textbf{Python} como el lenguaje de programación para el
desarrollo del análisis de datos en el Módulo~2. Dicha decisión 
satisface de forma estricta los requerimientos funcionales descritos 
en la \cref{it:RF04}, garantizando una integración nativa con el formato 
JSON utilizado por el Módulo 1 mediante su biblioteca estándar \texttt{json}, 
lo que simplifica el intercambio de información sin requerir dependencias 
externas complejas.

En relación con la complejidad computacional del modelo de 
Mamdani asociada principalmente al cálculo de la desfuzificación 
por centroide, la selección de Python como lenguaje de ejecución
resulta técnicamente idónea para el Módulo~2. Si bien Python es 
un lenguaje interpretado, bibliotecas especializadas como
\texttt{scikit-fuzzy} y \texttt{NumPy} delegan el cálculo matricial 
e integración numérica a rutinas internas. 
Esta integración vectorial reduce drásticamente 
los tiempos de procesamiento.

En contraste, se descartó C\# debido a su alta rigidez y dependencia 
del ecosistema .NET, lo que incrementa el consumo de memoria y dificulta 
la flexibilidad de despliegue. Por su parte, aunque R presenta fortalezas 
en el análisis estadístico matricial, su sintaxis especializada y su baja 
versatilidad para el desarrollo web e integración con microcontroladores
limitan su aplicabilidad en la arquitectura telemática propuesta.





\begin{landscape}
\begingroup
\setlength{\tabcolsep}{.1pt} % Espacio entre columnas reducido
\renewcommand{\arraystretch}{1.5} % Altura de las filas ligeramente reducida

% 1. El entorno longtable se reemplaza por tabular.
% 2. La tabla completa está dentro de \resizebox para forzar el ajuste al ancho total (landscape).
\begin{table}[h!]
\centering
\caption{Lenguajes de programación: C\#, Python y R}
\label{tab:lenguajesprogramacion}
\footnotesize 
% <-- Ajusta la tabla al ancho de la página horizontal
    % El comando \small reduce el tamaño de la fuente para ayudar al ajuste
    \small 
    \begin{tabular}{|>{\raggedright\arraybackslash}p{2.5cm}|>{\raggedright\arraybackslash}p{6.7cm}|>{\raggedright\arraybackslash}p{7cm}|>{\raggedright\arraybackslash}p{6.7cm}|}
    \hline
    \textbf{Aspecto} & \textbf{C\#}\cite{MarceloVillalobos2008}\cite{M.Deitel2007} & \textbf{Python}\cite{GonzalezDuque}\cite{Downey2002} & \textbf{R}\cite{Paradis2003} \\
    \hline
    \textbf{Caracterí- sticas} & Lenguaje de programación visual orientado a objetos, tipado, estático y controlado por eventos mediante el uso de un IDE. Desarrollado e integrado para el ecosistema .NET  & Lenguaje multiparadigma ya que soporta una programación orientada a objetos, imperativa y funcional con tipado dinámico y con una sintaxis muy fácil de aprender & Lenguaje especializado en análisis estadístico con un enfoque en manipulación de datos y generación de gráficos. \\
    \hline
    \textbf{Rendimi- ento} & Compilado en MSIL y luego optimizado mediante JIT en .NET, de nivel medio. & Interpretado, de alto nivel, lo que implica menor velocidad en tareas intensivas en CPU. Pero permite una gran portabilidad en diferentes entornos de programación. & Interpretado, de alto nivel, optimizado para operaciones estadísticas y matriciales, pero menos eficiente en tareas generales. \\
    \hline
    \textbf{Escalabi- lidad} & Compatible con JSON a través de bibliotecas siendo ampliamente eficiente en sistemas distribuidos y escalables como en ASP.NET. & Soporte nativo para JSON con la biblioteca json y escalable de forma nativa con diferentes \textit{frameworks} para aplicaciones web y con bases de datos modernas. & Compatible con JSON mediante paquetes y escalable en entornos estadísticos de \textit{big data}.\\
    \hline
    \textbf{Ventaja} & 
    Permite una gran integración y soporte multiplataforma robusta con .NET Core al tener un rendimiento optimizado y un tipado estático lo que reduce errores en tiempo de ejecución.
    & 
    Al tener una curva de aprendizaje suave su sintaxis es muy simple y clara, teniendo un ecosistema extenso y una comunidad activa lo que permite una amplia gama de bibliotecas para ciencia de datos y otras herramientas.
    & 
    Permite una visualización muy especializada y un análisis estadístico muy avanzado, siendo que tiene un gran soporte en sus bibliotecas especializadas en el análisis de datos respaldadas por una gran comunidad académica.
    \\
    \hline
    \textbf{Desventaja} & 
    Menor flexibilidad al tener mucha dependencia del ecosistema .NET Core lo que hace gran dependencia a este ecosistema. Esto también hace que algunas herramientas o bibliotecas puedan tener algún costo, aparte el ecosistema .NET puede ser algo robusto en cuestión de memoria.
    & 
    Tiende a tener un rendimiento limitado debido a una concurrencia y un tipado dinámico al hacer la interpretación de línea por línea, lo que lo hace más lento en aplicaciones intensivas. Haciendo que tenga dependencia de bibliotecas externas para algunas aplicaciones.
    & 
    Tiene una flexibilidad muy baja con otros entornos enfocados al desarrollo web, esto debido a su sintaxis única lo que hace que tenga una curva de aprendizaje más pronunciada.
    \\
    \hline
    \end{tabular}
\end{table}

\endgroup
\end{landscape}
